\section{Definitions}\label{sec:offline-definitions}

\subsection{System entities}

An offline payment system involves five types of participants: the
\textbf{central bank}, \textbf{intermediaries}, \textbf{users}, a
\textbf{registration authority}, and an \textbf{auditor}.

The central bank issues tokens to users and manages the reserves of the
intermediaries. Issuance is triggered by a \emph{withdrawal} request from an
intermediary and, when successful, debits the issued value from that
intermediary's reserves. The central bank also redeems tokens on a
\emph{deposit} request from an intermediary, crediting the redeemed value to
the requesting intermediary's reserves.

Users hold bank deposits at intermediaries, which they can exchange for
tokens. A user withdraws a token by submitting a withdrawal request to their
intermediary, which routes the request to the central bank; a successful
withdrawal debits the withdrawn value from the user's deposits and from the
intermediary's reserves. Once a user holds a token, they can use it in
offline payments. The token's lifecycle ends with a deposit, which lets its
owner claim bank deposits in exchange for the token; a deposit succeeds only
if the token was issued by the central bank, was initiated by its owner, and
had not previously been deposited.

Intermediaries maintain users' bank deposit accounts and mediate the
interaction of users with the central bank. They also onboard users into the
system, opening accounts for them and helping them obtain long-term
credentials from the registration authority. The registration authority
issues these long-term credentials, which users subsequently use to authorise
and authenticate payments. The auditor is authorised to trace a fraudulent
offline payment, such as one that double spends a token, to its originator.

We instantiate each user's wallet with a \textbf{mobile device} and a
\textbf{secure element}, which together let the user perform offline payments
securely and accountably.

We assume the registration authority, the central bank, and the
intermediaries are honest but curious: they follow the protocol, but may
collude with other participants to learn information about offline payments.
The registration authority is additionally trusted to assign each user a
unique identifier, to generate a correctly formed credential that will be
accepted by all other participants, and to issue that credential only
following well-established KYC processes. Users may be arbitrarily malicious:
they may collude to forge tokens, double spend them, or steal the tokens of
honest users, and may collude with other participants to undermine the
privacy of honest users. Users have full control of their mobile devices, but
we assume they cannot easily tamper with their secure elements. The auditor
is fully trusted to inspect only fraudulent payments and to de-anonymise only
misbehaving users. We consider static corruption throughout: the adversary
fixes the set of corrupt participants before any protocol interaction takes
place.

\subsection{Security goals}

An offline payment system must protect against forgery, token theft, and
double spending, and must preserve the privacy of users against the central
bank and the intermediaries. More specifically, it must ensure:

\begin{description}

  \item[Unforgeability.] An honest payee should not be tricked, during a
    payment, into accepting a token that was not issued by the central bank.
    Similarly, the honest-but-curious central bank will not declare a deposit
    successful unless it previously issued the corresponding token.

  \item[Theft prevention.] A malicious payer should not be able to convince
    an honest payee to accept a token the payer does not own. Likewise, a
    malicious user should not be able to successfully deposit a token they do
    not own.

  \item[Double-spending detection and identification.] An honest user uses a
    token only once, in either a payment or a deposit, after which the token
    is deleted; a malicious user may instead use a token in multiple payments
    or deposits. A secure offline payment system must ensure that such double
    spending is swiftly detected at the time of deposit, and that the
    malicious user is identified by the authorised auditor.

  \item[Non-frameability.] An honest user should never be accused of double
    spending by the authorised auditor, even if every other user, the
    intermediaries, and the central bank collude.

  \item[User anonymity.] Only the intermediary of an honest user making a
    withdrawal should learn her identity. During a payment, the history of
    the token being transferred leaks no information about the honest users
    involved in previous payments beyond what is unavoidably revealed:
    namely, that the payee of the last payment is the payer of the current
    one, and a growing list of random values fixed throughout the token's
    lifetime. We assume the payer and the payee know each other, and that
    this knowledge is not derived from the payment transcript. The same holds
    for deposits and withdrawals: the intermediary of a withdrawing or
    depositing user knows their identity, but the history of a deposited
    token leaks nothing about the honest users involved in payments
    preceding the deposit, beyond the list of random values associated with
    the token and the intermediary that identifies the depositor for
    accounting purposes.

  \item[Token unlinkability.] The intermediaries and the central bank should
    not be able to link the withdrawal of an honest user to a deposit, beyond
    what is revealed by the value of the withdrawal and the deposit, even if
    the intermediaries and the central bank collude with every other user in
    the system.

\end{description}

\subsection{Ideal functionality $\F$}\label{sec:ideal-func}

We define security for the offline payment protocol through an ideal
functionality $\F$, depicted in Figure~\ref{fig:ideal-func}, which represents
an incorruptible trusted party that receives the inputs of each protocol task
and returns the prescribed outputs. An adversary in the ideal world interacts
with $\F$ only through the parties it corrupts, which ensures that it can
only determine the inputs of corrupt parties and observe their outputs,
rather than making $\F$ deviate from its specification. Following
convention, we call this adversary a \emph{simulator} and denote it $\simu$.

We say that a protocol $\prot$ securely realises $\F$ if, for every adversary
$\adv$ that controls a subset of the parties and interacts with $\prot$,
there exists a simulator $\simu$ that controls the same parties and interacts
with $\F$, such that for every input of the corrupt parties, the output of
$\adv$ in the real world is indistinguishable from the output of $\simu$ in
the ideal world. Formally, $\prot$ securely realises $\F$ if and only if
$\real_{\adv, \prot} \approx \ideal_{\simu, \F}$, where $\real_{\adv, \prot}$
is the view of $\adv$ interacting with $\prot$ in the real world, and
$\ideal_{\simu, \F}$ is the view of $\simu$ interacting with $\F$ in the
ideal world.

\paragraph{Participants.} We identify the offline payment participants
$\usrs = \{\usr_1, \ldots, \usr_n\}$, who act as both payers and payees, the
intermediaries $\I = \{\Int_1, \ldots, \Int_m\}$, the auditor $\aud$, the
registration authority $\reg$, and the central bank $\cb$. For simplicity we
assume each user $\usr_i$ has one bank account and is equipped with a mobile
device $\dev_i$ and a secure element $\se_i$, so that $\usr_i$ is defined as
the pair $(\dev_i, \se_i)$; an individual with multiple bank accounts is
accommodated by mapping them to multiple users, each with one account. We
write $\SE = \{\se_1, \ldots, \se_n\}$ and $\Dev = \{\dev_1, \ldots, \dev_n\}$
for the sets of secure elements and mobile devices. As above, $\cb$, $\reg$,
and the intermediaries $\I$ are honest but curious, so the simulator $\simu$
accesses their inputs and outputs by default, and we assume a static
corruption model in which $\simu$ corrupts the parties of its choosing before
they interact with $\F$.

\paragraph{Data stores.} $\F$ maintains several data stores. $\reserves$
tracks each intermediary's reserves at $\cb$, with entries $\reserves[\Int] =
v$. $\registered$ is the set of registered users. $\accounts$ tracks user
accounts, with entries $\accounts[\usr] = \langle \Int, v \rangle$ indicating
that $\usr$ holds an account of value $v$ with intermediary $\Int$.
$\tokens$ tracks the history of token ownership, with entries
$\tokens[\vec{\rho}] = \langle \usr, v \rangle$ indicating that $\usr$ owns
the token of value $v$ whose history is identified by $\vec{\rho}$.
$\deposited$ tracks the history of deposited tokens, with entries
$\deposited[\rho] = \{\vec{\rho}_1, \ldots, \vec{\rho}_n\}$ such that $\rho$
identifies a withdrawn token and each $\vec{\rho}_i$ identifies a chain of
payments involving that token and ending in a deposit. Finally, $\mal$ is the
set of corrupted parties in $\Dev$ or $\SE$. All of $\registered$,
$\accounts$, $\tokens$, $\deposited$, and $\mal$ are initially empty.

\paragraph{Interfaces.} $\F$ provides the following interfaces to $\simu$ and
the participants.

\underline{$\corrupt$}. When $\simu$ wishes to compromise a party $P \in \Dev
\cup \SE$, it issues $\corrupt$, adding $P$ to $\mal$. Corrupting a user's
device without corrupting their secure element is possible; since
corrupting a secure element is much harder than corrupting a device, we
assume for simplicity that if $\se$ is corrupt so is the associated $\dev$.

\underline{$\fcreate$}. An intermediary $\Int$ calls $\fcreate$ to create an
account for user $\usr \in \usrs$ with initial value $v$, adding entry
$\accounts[\usr] \gets \langle \Int, v \rangle$ and updating $\Int$'s
reserves; $\F$ enforces that $\usr$ has one account in the system.

\underline{$\fregister$}. The registration authority $\reg$ calls
$\fregister$ to enrol a user, for whom an account was already created, in the
system. Users can be registered only once, and registration adds $\usr$ to
$\registered$.

\underline{$\fwithdraw$}. An intermediary $\Int$ indicates a value and
identifies a registered user $\usr$ whose account it hosts. If successful,
$\usr$ obtains a token identified by a unique random number $\rho$, usable
subsequently in offline payments authorised by $\usr$'s secure element; the
call succeeds only if $\usr$'s account and $\Int$'s reserves hold enough
value. To ensure a withdrawal is not linkable to a future deposit, the
identifier $\rho$ must not be known to $\Int$ or $\cb$: it is selected by
$\F$ if $\usr$'s secure element is honest, and by $\simu$ otherwise. Once
$\rho$ is selected and confirmed unique, $\F$ adds entry $\tokens[(\rho)]
\gets \langle \usr, v \rangle$. Since $\Int$ and $\cb$ are honest but
curious, $\simu$ learns the identity of the user and the value of the
withdrawal by default; if either the secure element or the mobile device of
the user is corrupt, $\simu$ additionally learns $\rho$.

\underline{$\fpay$}. A registered user $\usr$ invokes $\fpay$ to transfer
ownership of a token they own to some user $\usr'$, identifying the payer,
the payee, the value of the token, and a vector $\vec{\rho}$ that lets $\F$
locate the token in $\tokens$ and trace its history. $\F$ checks that $\usr'$
would like to accept the payment and that their mobile device and secure
element are not corrupt; if so, $\F$ checks that the token identified by
$\vec{\rho}$ is indeed owned by $\usr$ and has value $v$. If so, $\F$
produces a fresh random identifier $\rho'$ that, together with $\vec{\rho}$,
uniquely identifies the payment, and adds entry $\tokens[\vec{\rho} \|
\rho'] \gets \langle \usr', v \rangle$, letting a registered $\usr'$
subsequently transfer the token onward. If either the mobile device or the
secure element of $\usr'$ is corrupt, $\F$ does not condition success on the
validity of the token, reflecting that a malicious payee may choose to
accept an invalid one; in this case $\rho'$ is selected by $\simu$ if
$\usr'$'s secure element is corrupt, and by $\F$ otherwise, and
$\tokens[(\vec{\rho}, \rho')]$ is updated only if the token being
transferred is genuinely valid and $\usr'$ is registered. This ensures a
malicious payee who accepts an invalid token cannot transfer it onward to
honest users, since it will not appear in $\tokens$. If the mobile device or
secure element of either $\usr$ or $\usr'$ is corrupt, $\simu$ learns the
identities of both; otherwise, $\simu$ learns only the value and $\vec{\rho}
\| \rho'$. We stress that an unregistered payee can accept a payment and
receive a token, but can never transfer its ownership onward, as it is never
added to $\tokens$.

\underline{$\fdeposit$}. An intermediary $\Int$ invokes $\fdeposit$ to let a
registered user exchange a token for commercial bank money, identifying the
depositing user $\usr$, a vector $\vec{\rho}$ referring to the token and its
history, and the deposit value. $\F$ verifies that the token being deposited
is indeed owned by $\usr$ and has the right value, then updates
$\deposited[\vec{\rho}[0]] \gets \deposited[\vec{\rho}[0]] \cup
\{\vec{\rho}\}$ and checks whether $|\deposited[\vec{\rho}[0]]| > 1$; if so,
this signals double spending and $\F$ stops. Otherwise $\F$ updates the
holdings of $\usr$ and $\Int$. Since $\cb$ and $\Int$ are honest but
curious, $\simu$ learns the deposit information by default: the identity of
$\usr$, the value $v$, and the vector $\vec{\rho}$.

\underline{$\faudit$}. The auditor $\aud$ calls $\faudit$ with a token
identifier $\rho$, the identifier used during the corresponding withdrawal.
$\F$ retrieves $\deposited[\rho]$; if it contains nothing or a single
vector, $\F$ returns nothing, signalling that there is no double spender to
identify. Otherwise, $\F$ retrieves the set of vectors $\{\vec{\rho}_1,
\ldots, \vec{\rho}_l\}$ stored in $\deposited[\rho]$, and for each pair
$\vec{\rho}_i, \vec{\rho}_j$ identifies the double spender as the user
$\usr_{ij}$ stored in $\tokens[\vec{\rho}_{ij}]$, where $\vec{\rho}_{ij}$ is
the common prefix of $\vec{\rho}_i$ and $\vec{\rho}_j$. $\F$ returns the set
of such users to $\aud$ and to $\cb$.

\begin{figure*}[p]
  \fbox{\parbox{0.98\linewidth}{
    \centering
    \begin{algorithmic}[1]
      \scriptsize
      \begin{multicols}{2}

      \Statex Before any call to the offline payment interfaces, $\simu$
      issues corruption queries:
      \smallskip

      \Statex \underline{$\corrupt$} Upon an active corruption request
      $(\corrupt, P)$ for $P \in \Dev \cup \SE$ from $\simu$ do: $\mal \gets
      \mal \cup P$.
      \smallskip

      \Statex \underline{$\fcreate$} Upon receiving for the first time
      $(\fcreate, \usr, v)$ from $\Int \in \I$ for $\usr \in \usrs$, send
      $(\fcreate, \usr, \Int, v)$ to $\simu$ and wait for its go-ahead. On
      receipt of the go-ahead do:
      \begin{enumerate}
        \item {\bf If} $\accounts[\usr] \neq \perp$ {\bf then} stop
        \item {\bf Else}
        \begin{enumerate}
          \item $\accounts[\usr] \leftarrow \langle \Int, v\rangle$
          \item $\reserves[\Int] \leftarrow \reserves[\Int]+v$
          \item return $\fin$ to $\Int$
        \end{enumerate}
      \end{enumerate}
      \smallskip

      \Statex \underline{$\fregister$} Upon receiving for the first time
      $(\fregister, \usr)$ from $\reg$ for $\usr \in \usrs$ that already has
      an account -- send $(\fregister, \usr)$ to $\simu$ and wait for its
      go-ahead. On receipt of the go-ahead do:
      \begin{enumerate}
        \item $\registered \gets \registered \cup \usr$
        \item return $\fin$ to $\reg$
      \end{enumerate}
      \smallskip

      \Statex \underline{$\fwithdraw$} On receiving a $(\fwithdraw, \usr =
      (\se, \dev), \Int, v)$ from registered user $\usr$ that has an account
      with $\Int$, send $(\fwithdraw, \usr, \Int, v)$ to $\Int$ and
      $(\fwithdraw, \Int, v)$ to $\cb$ and await their go-ahead. Then send
      message $\fwithdraw$ to $\simu$ and await its go-ahead. On receipt of
      the go-ahead, fetch $\accounts[\usr] \rightarrow \langle \Int,
      v^*\rangle$, and do:
      \begin{enumerate}
        \item {\bf If} $v^* < v$ {\bf or } $\reserves[\Int] < v$ {\bf then}
        stop
        \item {\bf Else} do:
        \begin{enumerate}
          \item $\accounts[\usr] \leftarrow \langle \Int, v^* - v\rangle$
          \item $\reserves[\Int] \leftarrow \reserves[\Int] - v$
          \item {\bf If} $\se \not\in \mal$ {\bf then} randomly select $\rho$
          \item {\bf Else} receive $\rho$ from $\simu$
          \item {\bf If} $\tokens[(\rho)] \neq \perp$ {\bf then} stop
          \item {\bf Else} do:
          \begin{enumerate}
            \item $\tokens[(\rho)] \gets \langle \usr, v \rangle$
            \item {\bf If} $\se \in \mal$ {\bf or} $\dev \in \mal$ {\bf then}
            return $(\fwithdraw, \usr, \rho, v)$ to $\simu$
            \item {\bf Else} return $(\fwithdraw, \usr, v)$ to $\simu$
            \item return $(\fwithdraw, \rho, v, \fin)$ to $\usr$
          \end{enumerate}
        \end{enumerate}
      \end{enumerate}
      \smallskip

      \Statex \underline{$\fpay$} On receiving $(\fpay, \vec{\rho}, v, \usr' =
      (\se', \dev'))$ from registered user $\usr = (\se, \dev)$, send
      $(\fpay, \vec{\rho}, v)$ to $\usr'$. On receiving the response $\accept$
      from $\usr'$, send $(\fpay, \vec{\rho}, v)$ to $\simu$ and await the
      go-ahead. On receipt of the go-ahead do:
      \begin{enumerate}
        \item {\bf If} $\se' \not\in \mal$ {\bf and} $\dev' \not\in \mal$
        {\bf then} do:
        \begin{enumerate}
          \item {\bf If} $\tokens[\vec{\rho}] \neq \langle \usr, v \rangle$
          {\bf then} stop
          \item {\bf Else} do:
          \begin{enumerate}
            \item randomly select a unique random number $\rho'$
            \item $\tokens[\vec{\rho'}] \gets \langle\usr', v \rangle$ where
            $\vec{\rho'} = (\vec{\rho}, \rho')$
            \item {\bf If} $\se \in \mal$ {\bf or} $\dev \in \mal$ {\bf then}
            return $(\fpay, \vec{\rho'}, v, \usr, \usr')$ to $\simu$
            \item {\bf Else} return $(\fpay, \vec{\rho'}, v)$ to $\simu$
          \end{enumerate}
        \end{enumerate}
        \item {\bf Else} do:
        \begin{enumerate}
          \item {\bf If} $\se' \in \mal$ {\bf then} receive $\rho'$ from
          $\simu$
          \item {\bf Else} randomly select $\rho'$
          \item {\bf If} $\tokens[\vec{\rho}] = \langle \usr, v \rangle$
          {\bf and} $\usr' \in \registered$ {\bf then} $\tokens[\vec{\rho'}]
          \leftarrow \langle \usr', v \rangle$
          \item send $(\fpay,\vec{\rho'}, v, \usr, \usr')$ to $\simu$
        \end{enumerate}
        \item send $(\fpay, \vec{\rho'}, v , \fin)$ to $\usr'$
      \end{enumerate}
      \smallskip

      \Statex \underline{$\fdeposit$} On receiving $(\fdeposit, \vec{\rho}, v,
      \Int)$ from registered user $\usr$ that has an account with $\Int$,
      send $(\fdeposit, \vec{\rho}, v, \Int, \usr)$ to $\Int$ and
      $(\fdeposit, v, \Int)$ to $\cb$ and await their go-ahead. Then send
      $(\fdeposit, \vec{\rho}, v, \Int, \usr)$ to $\simu$ and await the
      go-ahead. On receipt of the go-ahead do:
      \begin{enumerate}
        \item {\bf If} $\tokens[\vec{\rho}] \neq \langle \usr, v \rangle$
        {\bf then} stop
        \item {\bf Else} do:
        \begin{enumerate}
          \item $\deposited[\vec{\rho}[0]] \gets \deposited[\vec{\rho}[0]]
          \cup \{\vec{\rho}\}$
          \item {\bf If} $|\deposited[\vec{\rho}[0]]| > 1$ {\bf then} stop
          \item {\bf Else} do:
          \begin{enumerate}
            \item $\accounts[\usr] \rightarrow \langle \Int, v^*\rangle$
            \item $\accounts[\usr] \gets \langle \Int, v^*+v \rangle$
            \item $\reserves[\Int] \gets \reserves[\Int] + v$
            \item send $(\fdeposit, \fin)$ to $\usr$
          \end{enumerate}
        \end{enumerate}
      \end{enumerate}
      \smallskip

      \Statex \underline{$\faudit$} On receiving $(\faudit, \rho)$ from
      $\aud$, send $(\faudit, \rho)$ to $\simu$ and await the go-ahead. On
      receipt of the go-ahead do:
      \begin{enumerate}
        \item {\bf If} $|\deposited[\vec{\rho}[0]]| \leq 1$ {\bf then} send
        $(\faudit, \fin)$ to $\aud$
        \item {\bf Else} do:
        \begin{enumerate}
          \item $\deposited[\rho] \rightarrow \{\vec{\rho}_1, ...,
          \vec{\rho}_n\}$
          \item $\cDS \leftarrow \emptyset$
          \item $\forall i\neq j $ do:
          \begin{enumerate}
            \item let $\vec{\rho}_{ij}$ denote the common prefix, i.e.\
            $\vec{\rho}_{ij} = (\rho_0, ..., \rho_{k-1})$ such that
            $\vec{\rho}_i[l] = \vec{\rho}_j[l]$ for all $l < k$ and
            $\vec{\rho_i}[k] \neq \vec{\rho_j}[k]$
            \item $\tokens[\vec{\rho}_{ij}] \rightarrow \langle v,
            \usr_{ij}\rangle$
            \item $\cDS \gets \cDS \cup \usr_{ij}$
          \end{enumerate}
          \item send $(\faudit, \cDS, \fin)$ to $\aud$ and $\cb$
        \end{enumerate}
      \end{enumerate}
      \end{multicols}
    \end{algorithmic}
  }}
  \caption{Ideal functionality $\F$ for offline payments.}
  \label{fig:ideal-func}
\end{figure*}
